% %%%%%%%%%%%%%%%%%%%%%%%%%%%%%%%%%%%%%%%%%%%%%%%%%%%%%%%%%%%%%%%%%%%%%
% %% Optims %%
% RTL Optimizations
% %%%%%%%%%%%%%%%%%%%%%%%%%%%%%%%%%%%%%%%%%%%%%%%%%%%%%%%%%%%%%%%%%%%%%

\section{RTL Optimizations for SpDCache}
\label{sec:optim:rtloptims}

\lettrine{T}{he} \acrshort{rtl} implementation of \acrlong{spdc} presented in Chapter~\ref{cha:arch} was essential to validate our performance models and to analyze performance based on matrix structure.
However, as already noted several times in the previous chapter, the \acrshort{rtl} implementation is missing some features and could be optimized for better practicality and performance.
In this section, we detail the major features and optimizations that are still to be implemented.

% %%%%%%%%%%%%%%%%%%%%%%%%%%%%%%%%%%%%%%%%%%%%%%%%%%%%%%%%%%%%%%%%%%%%%
\subsection{Floating-Point Support Implementation}
\label{sec:optim:rtloptims:fp}

The key focus of this thesis has been the structure on the non-zero elements in the matrix, however, the data-type of the matrices plays an important role in their storage footprint and the arithmetic complexity.
Indeed, in scientific computing, matrices are primarily stored in double precision (64-bits) floating-point.
In other applications such as DNNs or graph analytics, lower precision formats (\ex binary, 16-bits floating-point, 8-bits integer) are frequently used.
To simplify the design complexity of \acrlong{spdc}, we chose to work with 64-bits integer format.
This results in the same data-traffic as double floating-point, but simplifies the design, as the arithmetic operations (accumulations) can be done in a single cycle.
To deploy \acrlong{spdc} for high-performance computing applications, it must be extended to support the double data-type.

To support the accumulation of double format data between \textit{st1} and \textit{st2} in the \gls{pipeline} (Section~\ref{sec:arch:pipeline}), we could either deploy a single-cycle high-performance \acrshort{fpu} with the area trade-offs, or extend the \gls{pipeline} to accommodate two-cycle, for example, floating-point accumulation.
The latter approach better suits our \textit{accumulation table}, which does not stall the \gls{pipeline} and can parallelize element-wise accumulations across stored 64-byte lines.
Additionally, full floating-point support requires extending the AXI port to handle element-wise \acrfull{rmw} operations beyond its native integer capabilities~\cite{ARM2023_AXI}.

Floating-point support is critical for \acrlong{spdc} deployment.
Our expectation is that the associated \acrshort{rtl} changes, while time-consuming, should not present fundamental technical obstacles.

\begin{highlight}{Conclusion}{Red}
    \textbf{Floating-point support} is achievable through \gls{pipeline} extension and requires extending the AXI port for element-wise operations, presenting no fundamental implementation issues.
\end{highlight}

% %%%%%%%%%%%%%%%%%%%%%%%%%%%%%%%%%%%%%%%%%%%%%%%%%%%%%%%%%%%%%%%%%%%%%
\subsection{Region-Sizes Configuration}
\label{sec:optim:rtloptims:config}

One strength of \acrlong{spdc}'s \acrshort{rtl} implementation is its highly configurable design, inherited from the HPDcache implementation used as a foundation.
This allowed us to develop a single \acrshort{rtl} codebase that could be compiled and simulated for multiple designs, \acrshort{gc}, \acrshort{redc}, and \acrshort{spdc}, each with different region size configurations.
However, these configurations are static, requiring \acrshort{rtl} recompilation for each design variant.

Currently, \acrshort{spdc}'s \gls{reduction} mode can be dynamically enabled or disabled via an instruction issued before and after the \acrshort{spmv} kernel.
This switches between a generic cache and a hard-coded \acrshort{rtl} configuration.
To enhance software flexibility, \acrshort{spdc} should support dynamic configuration settings through the initial configuration instruction.

A fundamental limitation is that region sizes are restricted to powers of two, which severely constrains configuration options and led to the 25/50/25\% configuration used in the previous chapter.
This restriction stems from address decoding that uses bit shifts to extract the \gls{set} index, an approach inherited from HPDcache.
Supporting arbitrary region sizes would require integer division and modulo operations, which could introduce long, critical paths.
However, our Python simulations in Chapter~\ref{cha:spdc} demonstrate that precise region sizing is unnecessary for achieving good performance, making this optimization less important.

A practical compromise would be to provide one or two predefined configurations for \acrshort{redc} and \acrshort{spdc}, allowing dynamic selection by software while maintaining a simple user interface.

\begin{highlight}{Conclusion}{Red}
    Dynamic region-size configuration through software, while desirable, has a high hardware complexity without significant performance gains.
    A practical approach is to provide \textbf{predefined configurations for dynamic selection}, balancing flexibility with implementation simplicity.
\end{highlight}

% %%%%%%%%%%%%%%%%%%%%%%%%%%%%%%%%%%%%%%%%%%%%%%%%%%%%%%%%%%%%%%%%%%%%%
\subsection{Hardware Calculation of Region and Band-Width}
\label{sec:optim:rtloptims:regionInHW}

\acrlong{spdc} evaluation shows that depending on matrix characteristics, the \acrshort{op} \acrshort{spmv} kernel can be compute-bound, making region calculation a critical software bottleneck.
We addressed this by using a custom sparse format that converts latency costs into memory traffic costs.
To eliminate any latency or memory traffic overhead, we could compute the region argument of the \acrshort{red} instruction directly in hardware.
This approach sacrifices software flexibility but enables \acrshort{csc} format compatibility and simplifies the \acrshort{op} \acrshort{spmv} kernel to the version shown in Algorithm~\ref{src:arch:op_spmv}.
The required \acrlong{spdc} modifications are:
\begin{itemize}
    \item Configuration: The initial instruction enabling \acrlong{spdc}'s \gls{reduction} mode must accept a configuration argument specifying the word size in bytes ($s_w$), the output vector address ($c$), and the \acrshort{drc} size in bytes ($s_\mathit{DRC}$).
    The hardware could fetch the last parameter independently.
    When this instruction is issued, prior to each \acrshort{spmv}, the half-\gls{bw} is computed and retained until reconfigured, where $s_l$ denotes cache-line size in bytes:
    $$w_\mathit{hB} = \dfrac{s_\mathit{DRC}-s_l}{2 \times s_w}+1$$
    The half-\gls{bw} could also be computed in software, and given as argument to the configuration.
    %
    \item Column indexing: Each outer-loop iteration requires issuing a new instruction containing the current column index $j \in \ldb 0; M-1 \rdb$, denoted $CCOL(j)$.
    This value persists until the next $CCOL$ instruction arrives, consuming one cycle per outer-loop iteration.
    Note that the \acrshort{reg} instruction becomes unnecessary.
    %
    \item Region determination: For each $RED(addr, val)$ instruction, \acrshort{spdc} computes the current row $i \in \ldb 0; M-1 \rdb$ from the address: $i = \dfrac{addr-c}{s_w}$, then determines the region as \acrshort{ibr} if $|i-j| \leqslant w_\mathit{hB}$, otherwise \acrshort{obr}.
\end{itemize}
These modifications impose minimal hardware overhead, with per-\gls{reduction} computations executed at \textit{st0} in the \acrlong{spdc} \gls{pipeline}.
This feature becomes essential when extending \acrshort{spdc} to multi-level, multi-core architectures (Section~\ref{sec:optim:multicore}).

\begin{highlight}{Conclusion}{Red}
    \textbf{Hardware-based region computation} eliminates unnecessary instructions in the software's inner-loop, enabling \acrshort{csc} format compatibility and simplifying the \acrshort{spmv} kernel with minimal hardware overhead.
\end{highlight}

% %%%%%%%%%%%%%%%%%%%%%%%%%%%%%%%%%%%%%%%%%%%%%%%%%%%%%%%%%%%%%%%%%%%%%
\subsection{Complete Inter-Region Cache Coherency}
\label{sec:optim:rtloptims:coherency}

When executing an \acrshort{op} \acrshort{spmv} kernel with \acrlong{spdc}, we assume that the processor issues \acrshort{red} instructions strictly on addresses of the output vector $c$, and that loads and stores are issued only on addresses not in the output vector $c$.
\acrlong{spdc}'s current implementation provides no hardware safeguards if this assumption is violated.
While we provide local coherency between the two \gls{reduction} regions (\acrshort{drc} and \acrshort{srt}) to prevent $c$ address duplication, we lack coherency with the \acrshort{grc}.
A simple detection system could verify that \glspl{reduction} are issued exclusively on addresses within the output vector's range, while other instructions operate outside this range.
PHI~\cite{Mukkara2019_PHI} demonstrates that full coherency is achievable, as required in their design, proving that \glspl{reduction} and generic instructions can safely coexist within the same cache hierarchy.
In \acrlong{spdc}'s case, implementing global inter-region coherency would require the following modifications:
\begin{itemize}
    \item Configuration: When enabling \gls{reduction} mode, the configuration parameter must specify the address range for \acrshort{red} instructions (\eg, output vector $c$ and its extent).
    %
    \item Region Assignment: To prevent address duplication between \gls{reduction} regions and the \acrshort{grc}, loads and stores on $c$ addresses target \gls{reduction} regions, while \acrshortpl{red} on non-$c$ addresses target the \acrshort{grc}.
    The following cases apply only when \gls{reduction} mode is enabled.
    %
    \item Load on $c$ address:
    \begin{itemize}
        \item \Glspl{miss} in \acrshort{drc} and \acrshort{srt}: Issue memory read and return data to processor, bypassing cache.
        \item \Gls{hit} in \acrshort{drc} or \acrshort{srt}: Issue memory read, accumulate with cached data, and return result to processor.
    \end{itemize}
    We do not update \gls{reduction} regions on load to avoid complexity.
    Updating would require writing zeros to memory to prevent double accumulation upon \gls{eviction} or flush.
    %
    \item Store on $c$ address:
    \begin{itemize}
        \item \Glspl{miss} in \acrshort{drc} and \acrshort{srt}: Issue memory write with store value, bypassing cache.
        \item \Gls{hit} in \acrshort{drc} or \acrshort{srt}: Overwrite cached data with store value and issue memory write of a zero to maintain consistency.
    \end{itemize}
    Cache policy (write-through and write-back) can influence this design choice.
    Crucially, data must not exist in both memory and cache simultaneously.
    %
    \item \acrshort{red} on non-$c$ address:
    \begin{itemize}
        \item \Gls{miss} in \acrshort{grc}: Issue memory read, accumulate with \acrshort{red} value, insert into \acrshort{grc} (with \gls{eviction} if needed), and mark dirty (write-back) or write to memory (write-through).
        \item \Gls{hit} in \acrshort{grc}: Accumulate cached data with \acrshort{red} value, update \acrshort{grc} accordingly, and mark dirty (write-back) or write to memory (write-through).
    \end{itemize}
\end{itemize}
It is to be expected that incorrect region usage increases memory traffic and latency, degrading performance.
Implementing a coherency or detection scheme would ensure deterministic behavior, a clear benefit for \acrlong{spdc} deployment.

\begin{highlight}{Conclusion}{Red}
    \textbf{Complete inter-region cache coherency} prevents address duplication across \gls{reduction} and generic regions, ensuring deterministic behavior at minimal hardware cost through address-range configuration and region-specific load/store/reduction handling.
\end{highlight}